\subsection{The complete exclusion through twelve arrows}

\begin{theorem}\label{C:CP2-lower13}
Every finite category whose classifying space has the homotopy
type of $\mathbb CP^2$ has at least thirteen arrows.
Consequently the category minimum for $\mathbb CP^2$ lies in
the interval $[13,40]$ with the explicit upper model below.
\end{theorem}
\begin{proof}
Proposition~\ref{C:CP2-lower12} already excludes at most eleven
arrows. We therefore consider exactly twelve. A connected
category with $k$ objects has its $k$ identities and at least
$k-1$ further arrows connecting its underlying undirected
object graph. Thus $12\geq2k-1$ and $k\leq6$; there are no
additional cases with seven or more objects.

For one object the stronger monoid lower bound in
Corollary~\ref{B:fourdim-lower20} applies.
For two objects, categories with arrows in just one direction
are excluded by Lemma~\ref{C:twelve-triangular}, and those
with arrows in both directions by
Proposition~\ref{C:twelve-strongly-connected}.
Disconnected categories cannot have the connected homotopy
type under consideration.

For three objects,
Proposition~\ref{C:twelve-three-condensation} reduces the
problem to three singleton strongly connected components in
a linear order. The program
\path{code/c_categories12_linear.cpp} enumerates precisely
this remaining scope. Label the objects uniquely in increasing
order $0<1<2$; each of the three upper cross hom-sets is
nonempty and every lower cross hom-set is empty. Identity
arrows are labelled $0,1,2$, and the other arrows are sorted
by their source and target. The program enumerates every
such hom-cardinality profile and every type-correct product,
propagating all associativity equations. It discards only the
proved singleton cross-row or cross-column joins of
Lemma~\ref{C:singleton-cross-join}, and the previously proved
cofinal-object deletions. There is no comparison with an
incomplete earlier enumeration.

For four, five, and six objects, the complete
\path{code/c_categories12_cp2.cpp} runs enumerate all
connected hom-cardinality profiles and all associative
type-correct compositions after discarding initial- or
terminal-object profiles. All four completed enumeration
scopes have the following exact counts. The table count is
the number of associative labelled tables in the retained
scope; candidates have $H^1(-;\mathbb F_2)=0$ and
$\dim H^2(-;\mathbb F_2)=1$. The class count identifies
isomorphic or opposite candidate categories.
\[
\begin{array}{c|r|r|r|r|r}
 k&\text{tables}&\text{candidates}&\text{classes}
      &x^2=0&x^2\ne0,\ H^3\ne0\\\hline
 3\text{, linear}&1\,205\,250&165\,629&3\,624&3\,500&124\\
 4&36\,993\,658&3\,094\,704&4\,554&4\,545&9\\
 5&718\,340&77\,260&220&220&0\\
 6&51\,984&6\,000&6&6&0
\end{array}                                                    \tag{\ref{C:CP2-lower13}.1}
\]

Here is the independently checked obstruction for every
candidate class. The Python routine
\path{code/a_category_cups.py} verifies the identity,
source-target and associativity equations anew, constructs
the complete normalized nerve in the needed degrees, and
uses exact binary elimination to find the nonzero degree-two
cohomology class $x$. It computes its Alexander--Whitney
square. Every asserted zero square has an explicit degree-three
cochain primitive, checked on every normalized four-simplex.
For the 133 classes with nonzero square, a separate full
boundary-rank calculation gives $H^3(-;\mathbb F_2)\ne0$.
The generator's early $H^1,H^2$ tests and all these later
tests are necessary conditions for $\mathbb CP^2$, so each
case is excluded without inferring an all-degree homotopy
type from a homology cutoff.

The three-object stream is checked by
\path{code/c_twelve_stream_cups.py}, which marks completion
only after the generator's final summary exists and the
candidate file has been fully read. Its completed output is
\path{code/outputs/c_categories_12_3_linear_CP2_cups.json}:
all 165629 candidates have been read, all 3624 classes checked,
and the unexcluded list is empty.
The other completed outputs are
\path{code/outputs/c_categories_12_4_CP2_cups.json},
\path{code/outputs/c_categories_12_5_CP2_cups.json}, and
\path{code/outputs/c_categories_12_6_CP2_cups.json}.
The four-object file also contains its nine independent
residual rank calculations. Exact isomorphism caches use full
relabelled table strings, so a hash collision cannot suppress
a class. The known nonzero-square six-element class 190 is
rechecked as a positive control.

The earlier twelve-arrow streams labelled minimal, reduced,
or cofinal remain explicitly incomplete prefixes. None is
used for the completeness assertion in this theorem.
The completed scopes in (1), together with the structural
one-, two-, and three-object exclusions, exhaust every
possible object count and finish the proof.
\end{proof}
